\section{The special-flow history graph in full detail}
\label{subsec:invariant-history-details}

This subsection isolates the invariant-history part of
\cref{lem:shared-graph-trace}.  It does not address the two one-sided trace
problems.  The purpose is to specify the Banach spaces, the prepared
nonlinearities, the contraction, and the finite mixed prolongation without
appealing to a same-regularity implicit-function theorem.

The preparation in \cref{def:canonical-preparation} contains a tame-flow
condition in addition to its global derivative majorant.  We state that
condition in operator form below and verify it for the explicit cutoff in
the \((\chi,d)\)-coordinates.  The distinction is essential:
the global derivative bound alone would yield a factor
\(\delta A_\delta\exp(\Theta_*A_\delta)\), which need not tend to zero when
\(A_\delta=P(R_\delta)e^{cR_\delta}\).

\subsection{Prepared normal form and fixed history observations}

For the exact algebra in the first display, write
\(U_{XY}=\R^2\), with elements \(u=(X,Y)\), and retain the stable space
\(E_N=\ker\pi_N^T\) with the norm in \cref{eq:shared-norm}.  For current and
delayed states \(x_k=(u_k,h_k)\), \(0\leq k\leq m\), set
\begin{align}
 W_N(u_k,h_k)&=g_NX_k^2+h_k,\notag\\
 \mathscr L^X_{N,\zeta}[x_0,\ldots,x_m]
 &=\sum_{k=0}^m(B_{k,N}+\zeta R_{k,N})
       \mathbf1(X_0-X_k),\notag\\
 \mathscr L^W_{N,\zeta}[x_0,\ldots,x_m]
 &=\sum_{k=0}^m(B_{k,N}+\zeta R_{k,N})
       \{W_N(u_0,h_0)-W_N(u_k,h_k)\}.                 \label{eq:invdet-delay-polynomials}
\end{align}
The unprepared nonlinearities obtained from
\cref{eq:shared-exact-X,eq:shared-exact-h}, after eliminating the occurrence
of \(X'\) in the stable equation, are
\begin{align}
 F^{X,\natural}_N={}&-2X_0\pi_N^T\diag(c_N)W_N(u_0,h_0)
 -\frac\beta3X_0^3+K\pi_N^T\mathscr L^X_{N,\zeta}\notag\\
 &+\rho\{-\pi_N^T\diag(c_N)W_N(u_0,h_0)^{\circ2}
             +K\pi_N^T\mathscr L^W_{N,\zeta}\}\notag\\
 &-\rho^2\beta X_0\pi_N^TW_N(u_0,h_0)^{\circ2}
 -\frac{\rho^3\beta}{3}\pi_N^TW_N(u_0,h_0)^{\circ3},\notag\\
 F^{Y,\natural}_N={}&\nu,                              \label{eq:invdet-raw-F}\\
 G_N^\natural={}&P_{\perp,N}\{-2X_0\diag(c_N)W_N(u_0,h_0)
                    +K\mathscr L^X_{N,\zeta}\}\notag\\
 &-2g_NX_0\{Y_0-\alpha X_0^2+\rho F_N^{X,\natural}\}\notag\\
 &+\rho P_{\perp,N}\{-\diag(c_N)W_N(u_0,h_0)^{\circ2}
                  -\beta X_0^2W_N(u_0,h_0)
                  +K\mathscr L^W_{N,\zeta}\}\notag\\
 &-\rho^2\beta X_0P_{\perp,N}W_N(u_0,h_0)^{\circ2}
 -\frac{\rho^3\beta}{3}P_{\perp,N}W_N(u_0,h_0)^{\circ3}.       \label{eq:invdet-raw-G}
\end{align}
Thus, at \(\rho=\delta>0\), the exact chart has the form
\begin{equation}
 u'=q_0(u)+\rho F_N^\natural,
 \qquad
 \rho h'=A_Nh+\rho G_N^\natural.                     \label{eq:invdet-raw-normal-form}
\end{equation}
The factors \(\rho\) in both equations are exact, not formal Taylor
factors.

The graph transform itself is carried out in the \((\chi,d)\)-coordinates,
not in the \((X,Y)\)-norm.  Let
\[
 \mathcal K(X,Y)=
 \left(-2\alpha X,\,
       Y-\alpha X^2+\frac1{2\alpha}\right)=(\chi,d),
\]
put \(U=\R^2_{\chi d}\), \(X_N=U\times E_N\), and use the maximum product
norm.  Conjugating \cref{eq:invdet-raw-normal-form} by \(\mathcal K\)
preserves its exact form
\[
 u'=q_0^{\chi d}(u)+\rho\widehat F_N^\natural,\qquad
 \rho h'=A_Nh+\rho\widehat G_N^\natural,
 \qquad
 q_0^{\chi d}(\chi,d)=(1-2\alpha d,\chi d).
\]
Indeed, the difference between the transformed full reduced field and
\(D\mathcal K\,q_0\) still contains the external factor \(\rho\).
Derivatives of \(\mathcal K\) and \(\mathcal K^{-1}\) have only polynomial
growth on the logarithmic tube and are therefore absorbed in the tame
polynomial factors below.  From this point through the graph construction,
\(u\), \(Q\), and their norms refer to the \((\chi,d)\)-coordinates.
The final pointwise estimates are pulled back to \((X,Y)\).

Fix a target \(\delta\), freeze the graph preparation from
\cref{def:canonical-preparation}, and denote the resulting globally bounded
maps by
\begin{align}
 q_{0,\delta}&:U\longrightarrow U,\notag\\
 F_{N,\delta}&:X_N^{m+2}\times[-\delta,\delta]
          \times I_\nu\times[-\zeta_*,\zeta_*]\longrightarrow U,\notag\\
 G_{N,\delta}&:X_N^{m+2}\times[-\delta,\delta]
          \times I_\nu\times[-\zeta_*,\zeta_*]\longrightarrow E_N.
                                                               \label{eq:invdet-prepared-data}
\end{align}
They agree with the \(\mathcal K\)-conjugates of
\cref{eq:invdet-raw-F,eq:invdet-raw-G} on the uncut set.
The products in \cref{eq:shared-products}, the layer bound in
\cref{eq:shared-layer-bounds}, and the fixed number of atoms imply that all
the derivatives listed in \cref{eq:preparation-derivative-list} are bounded
in the displayed product norms without a factor depending on \(N\).  Write
\begin{equation}
 A_\delta=C_A\langle R_\delta\rangle^{M_A}e^{c_AR_\delta}          \label{eq:invdet-prepared-majorant}
\end{equation}
for one common majorant.  Since
\(R_\delta=O(\sqrt{\log(1/\delta)})\), for every fixed \(a>0\) and
\(k\in\mathbb N\),
\begin{equation}
 \delta^a A_\delta^k\longrightarrow0.                           \label{eq:invdet-subpower}
\end{equation}

The history observation is kept independent of \(\zeta\).  For a complete
vector field \(Q\) and a map \(H:U\to E_N\), define
\begin{align}
 \mathcal I_{Q,H}(u)(\vartheta)
 &=\bigl(\Phi_Q^\vartheta u,H(\Phi_Q^\vartheta u)\bigr),
       &&-\Theta_*\leq\vartheta\leq0,\label{eq:invdet-history-lift}\\
 \mathcal E_{Q,H}(u)
 &=\bigl(\mathcal I_{Q,H}(u)(0),
          \mathcal I_{Q,H}(u)(-\theta_0),\ldots,
          \mathcal I_{Q,H}(u)(-\theta_m)\bigr).       \label{eq:invdet-history-evaluation}
\end{align}
This is the atomic operator \(\mathbf M_N\) of
\cref{eq:shared-evaluation-TV} written as a tuple.  Its norm is at most
\(m+2\), uniformly in \(N\).  All \(\zeta\)-dependence remains in the
matrices in \cref{eq:invdet-delay-polynomials}; consequently no moving
Dirac mass is differentiated.  Moreover, the delay expressions are
differences, so the unbounded affine anchor \(u\mapsto(u,0)\) cancels.

The following tame-flow property is the exact input used below.

\begin{definition}[Tame frozen preparation]
\label[definition]{def:invdet-tame-preparation}
The frozen preparation is called tame through grade eleven if there are a
polynomial \(P_*\), constants \(c_*,\Gamma_*>0\), and an integer
\(M_*\geq1\), independent of \(N\) and \(\delta\), such that the order-zero
displacement estimate
\begin{equation}
 \|\Phi_Q^{\vartheta-\rho r}-\Id\|_{C_b(U;U)}
 \leq \mathfrak B_\delta(1+r^{M_*})
       e^{\Gamma_*(1+R_\delta)|\rho|r}               \label{eq:invdet-tame-displacement}
\end{equation}
holds for every vector field \(Q\) in the graph ball defined in
\cref{eq:invdet-graph-ball}.  At positive grade, it is required
inductively for every smooth Picard iterate whose preceding graph jets lie
in the invariant balls already constructed that
\begin{equation}
 \|\mathscr J\Phi_Q^{\vartheta-\rho r}\|_
 {C_b(U;\mathcal L_{\rm multi})}
 \leq \mathfrak B_\delta(1+r^{M_*})
       e^{\Gamma_*(1+R_\delta)|\rho|r},               \label{eq:invdet-tame-flow}
\end{equation}
for \(-\Theta_*\leq\vartheta\leq0\), \(r\geq0\),
\(|\rho|\leq\delta\), and every state--parameter flow jet
\(\mathscr J\) of total grade between one and eleven.  For two Picard
iterates in the corresponding common jet fibers, the difference of the
two flows (grade zero) and the difference of every positive-grade flow jet
obey the same right-hand side multiplied by the current-block difference
plus preceding-block differences, where
\begin{equation}
 \mathfrak B_\delta=P_*(R_\delta)e^{c_*R_\delta}.       \label{eq:invdet-tame-envelope}
\end{equation}
\end{definition}

This property is obtained by a concrete admissible choice.  Let
\(\omega\in C_c^\infty(\R)\) equal one on \([-1,1]\) and vanish off
\((-2,2)\).  Choose fixed \(B_*>2\Theta_*+B+3\) and \(d_*>0\) so that,
for all sufficiently small \(\delta\), the fixed \(b_0\)-neighborhood of
the depth-two hull in \cref{def:canonical-preparation} is contained in
\[
 |\chi|<R_\delta+B_*,\qquad |d|<d_*.
\]
Such a \(d_*\) exists because the normal multiplier over a fixed
depth-two singular-flow interval is \(e^{O(R_\delta)}\), while
\(\delta^\kappa e^{O(R_\delta)}=o(1)\).  In the
\((\chi,d)\)-coordinates take
\begin{equation}
 \Xi_\delta(\chi,d)
 =\omega\!\left(\frac{\chi}{R_\delta+B_*}\right)
  \omega\!\left(\frac d{d_*}\right),
 \qquad
 q_{0,\delta}^{\chi d}
 =\Xi_\delta(\chi,d)\binom{1-2\alpha d}{\chi d}.       \label{eq:invdet-explicit-q0-preparation}
\end{equation}
The Lipschitz estimate below is in the \((\chi,d)\)-norm:
\begin{equation}
 \operatorname{Lip}(q_{0,\delta})\leq C(1+R_\delta).  \label{eq:invdet-q0-lipschitz}
\end{equation}
Indeed, on the support of \(\Xi_\delta\),
\[
 |\chi|\leq2(R_\delta+B_*),\qquad |d|\leq2d_*,
\]
and hence
\[
 |q_0^{\chi d}|\leq C(1+R_\delta),\qquad
 \|Dq_0^{\chi d}\|
 =\left\|\begin{pmatrix}0&-2\alpha\\ d&\chi\end{pmatrix}\right\|
 \leq C(1+R_\delta).
\]
Moreover,
\[
 |\partial_\chi\Xi_\delta|\leq C(1+R_\delta)^{-1},
 \qquad
 |\partial_d\Xi_\delta|\leq C.
\]
The product rule therefore gives
\[
 \|D(\Xi_\delta q_0^{\chi d})\|
 \leq \|D\Xi_\delta\|\,|q_0^{\chi d}|
       +\|Dq_0^{\chi d}\|
 \leq C(1+R_\delta),
\]
which proves \cref{eq:invdet-q0-lipschitz}.  Higher derivatives are bounded
by polynomials in \(R_\delta\).  This estimate is not asserted for the
conjugated field in the Euclidean \((X,Y)\)-norm; converting back on the
retained tube costs polynomial factors in \(R_\delta\), already allowed in
\cref{eq:invdet-tame-envelope}.

For \(Q\) in the graph ball, its defining Lipschitz constraint
and \cref{eq:invdet-subpower} give
\begin{equation}
 \operatorname{Lip}(Q)
 \leq C(1+R_\delta)+b_{\rm gr,1}\delta\mathfrak B_\delta
 \leq C'(1+R_\delta).                                 \label{eq:invdet-Q-lipschitz}
\end{equation}
Since every prepared \(Q\) is bounded, integrating
\(\partial_t\Phi_Q^t=Q(\Phi_Q^t)\) proves
\cref{eq:invdet-tame-displacement}; no bound on the unbounded map
\(u\mapsto\Phi_Q^tu\) itself is asserted.  The ordinary flow variational
equations, successively differentiated through grade eleven in the block
order used below, then prove \cref{eq:invdet-tame-flow} for positive-grade
jets.  The state derivative of the perturbation in
\(Q\) retains \(\rho\), whereas a \(\rho\)-parameter derivative which
removes that factor is bounded by \(A_\delta\) and belongs to the prescribed
tame envelope.  Fixed-time flow factors are bounded by
\(P(R_\delta)e^{cR_\delta}\), and substitution of the terminal time
\(\vartheta-\rho r\) introduces only powers of \(r\) and the last
exponential in \cref{eq:invdet-tame-flow}.  This is a joint induction with
the invariant jet balls, not a smoothness assertion for arbitrary \(C^1\)
fields.  Thus
\cref{eq:invdet-explicit-q0-preparation}, together with the slotwise
preparation already required in \cref{def:canonical-preparation}, is a
sufficient concrete realization of
\cref{def:invdet-tame-preparation}.

\begin{remark}[Why the tame-flow clause is explicit]
\label[remark]{rem:invdet-missing-flow-hypothesis}
The bound in \cref{eq:preparation-derivative-list} permits
\(\operatorname{Lip}(q_{0,\delta})\) to be as large as \(A_\delta\).
Gronwall would then give \(e^{\Theta_*A_\delta}\) for a fixed delay
evaluation, and \cref{eq:invdet-subpower} does not control
\(\delta A_\delta e^{\Theta_*A_\delta}\).  This is why
\cref{def:canonical-preparation} includes the explicit singular cutoff and
requires \cref{eq:invdet-tame-flow}; its derivative majorant is not being
used as a substitute for flow control.
\end{remark}

\subsection{The order-zero graph transform}

For one target \(\delta\), put
\[
 \Lambda_\delta=[-\delta,\delta]_\rho\times I_\nu
                    \times[-\zeta_*,\zeta_*]
\]
and define
\[
 \mathfrak Z_N^0=C_b(U;U)\times C_b(U;E_N),
 \qquad
 \|(Q,H)\|_0=\max\{\|Q\|_\infty,\|H\|_\infty\}.
\]
The parameter-dependent pairs belong to
\(C(\Lambda_\delta;\mathfrak Z_N^0)\).  For fixed positive constants
\(b_{\rm gr,0},b_{\rm gr,1}\), to be chosen below, let
\begin{equation}
\begin{split}
 \mathbb B_{N,\delta}=\{(Q,H):{}&
 \sup_{\Lambda_\delta}\|(Q-q_{0,\delta},H)\|_0
       \leq b_{\rm gr,0}\delta A_\delta,\\
 &\sup_{\Lambda_\delta}
   \max\{\operatorname{Lip}(Q-q_{0,\delta}),
          \operatorname{Lip}(H)\}
       \leq b_{\rm gr,1}\delta\mathfrak B_\delta\}.
                                                               \label{eq:invdet-graph-ball}
\end{split}
\end{equation}
It is closed, hence complete, in the uniform metric.  Every \(Q\) in this
ball is globally bounded and Lipschitz and therefore has a complete
two-sided flow.

Define the special-flow transform by
\begin{align}
 \mathcal T_{Q,N}(Q,H)(u)
 &=q_{0,\delta}(u)+\rho F_{N,\delta}
       (\mathcal E_{Q,H}(u);\rho,\nu,\zeta),
                                                               \label{eq:invdet-TQ}\\
 \mathcal T_{H,N}(Q,H)(u)
 &=\rho\int_0^\infty e^{A_Nr}G_{N,\delta}
       (\mathcal E_{Q,H}(\Phi_Q^{-\rho r}u);
            \rho,\nu,\zeta)\,\dd r.               \label{eq:invdet-TH}
\end{align}
For \(\rho<0\), this is only an auxiliary fixed-point equation for Taylor
expansion.  Its RFDE interpretation is used only at \(\rho=\delta>0\).

\begin{lemma}[Self-mapping and contraction]
\label[lemma]{lem:invdet-C0-contraction}
Assume the frozen preparation is tame through grade eleven.  There is
\(\delta_0>0\), independent of \(N\), such that for
\(0<\delta<\delta_0\) the transform
\cref{eq:invdet-TQ,eq:invdet-TH} maps
\(\mathbb B_{N,\delta}\) into itself and is a contraction in the uniform
metric, with contraction factor at most \(1/2\).
\end{lemma}

\begin{proof}
Let \(Z=(Q,H)\) and \(\widetilde Z=(\widetilde Q,\widetilde H)\) be in the
ball.  The first flow variation with respect to the vector field gives,
for a fixed time \(t\),
\begin{equation}
 \sup_u|\Phi_Q^tu-\Phi_{\widetilde Q}^tu|
 \leq C\mathfrak B_\delta(1+|t|)e^{C(1+R_\delta)|t|}
       \|Q-\widetilde Q\|_\infty.                    \label{eq:invdet-flow-difference}
\end{equation}
This is the grade-zero two-flow difference clause in
\cref{def:invdet-tame-preparation}.  Combining it with the Lipschitz bound for
\(H\), and using the maximum norm on the fixed tuple in
\cref{eq:invdet-history-evaluation}, yields
\begin{align}
 \|\mathcal E_Z-\mathcal E_{\widetilde Z}\|_\infty
 &\leq C\mathfrak B_\delta\|Z-\widetilde Z\|_0,
                                                               \label{eq:invdet-E-difference}\\
 \|\mathcal E_Z(\Phi_Q^{-\rho r}\cdot)
       -\mathcal E_{\widetilde Z}
          (\Phi_{\widetilde Q}^{-\rho r}\cdot)\|_\infty
 &\leq C\mathfrak B_\delta(1+r^{M_*})
       e^{\Gamma_*(1+R_\delta)|\rho|r}
       \|Z-\widetilde Z\|_0.                         \label{eq:invdet-shifted-E-difference}
\end{align}
All constants are independent of \(N\); the only delay factor is the fixed
number \(m+2\).

Let \(a_*=D\gamma\).  After increasing the polynomial and exponential
constants in \(\mathfrak B_\delta\), the Lipschitz bounds for the prepared
nonlinearities and \cref{eq:poisson-bound} give
\begin{align}
 \|\mathcal T_{Q,N}(Z)-\mathcal T_{Q,N}(\widetilde Z)\|_\infty
 &\leq C\delta\mathfrak B_\delta\|Z-\widetilde Z\|_0,
                                                               \label{eq:invdet-TQ-difference}\\
 \|\mathcal T_{H,N}(Z)-\mathcal T_{H,N}(\widetilde Z)\|_\infty
 &\leq C\delta\mathfrak B_\delta
 \int_0^\infty(1+r^{M_*})
 e^{-\{a_*-\Gamma_*(1+R_\delta)\delta\}r}\,\dd r
 \|Z-\widetilde Z\|_0.                              \label{eq:invdet-TH-difference}
\end{align}
By \cref{eq:invdet-subpower}, after decreasing \(\delta_0\),
\begin{equation}
 \Gamma_*(1+R_\delta)\delta\leq a_*/2,
 \qquad
 q_\delta:=C\delta\mathfrak B_\delta
 \left(1+\int_0^\infty(1+r^{M_*})e^{-a_*r/2}\,\dd r\right)
 \leq\frac12.                                        \label{eq:invdet-contraction-factor}
\end{equation}
This proves contraction.

For self-mapping, boundedness of the prepared data gives
\begin{equation}
 \|\mathcal T_{Q,N}(Z)-q_{0,\delta}\|_\infty
 \leq\delta A_\delta,
 \qquad
 \|\mathcal T_{H,N}(Z)\|_\infty
 \leq\frac{\delta A_\delta}{a_*}.                    \label{eq:invdet-transform-size}
\end{equation}
The state version of \cref{eq:invdet-tame-flow}, with
\cref{eq:poisson-bound}, similarly gives
\begin{equation}
 \max\{\operatorname{Lip}(\mathcal T_{Q,N}-q_{0,\delta}),
          \operatorname{Lip}(\mathcal T_{H,N})\}
 \leq C\delta\mathfrak B_\delta.                    \label{eq:invdet-transform-Lipschitz}
\end{equation}
Choose \(b_{\rm gr,0}>\max\{1,a_*^{-1}\}\), then
\(b_{\rm gr,1}>C\), and decrease
\(\delta_0\) once more if necessary.  Equations
\cref{eq:invdet-transform-size,eq:invdet-transform-Lipschitz} make the ball
invariant.
\end{proof}

Banach's theorem now gives a unique fixed point
\begin{equation}
 Z_{N,\rho,\nu,\zeta}=(Q_{N,\rho,\nu,\zeta},
                       H_{N,\rho,\nu,\zeta})          \label{eq:invdet-fixed-graph}
\end{equation}
in the declared ball, simultaneously and continuously on
\(\Lambda_\delta\).  Uniqueness is asserted in this ball, not among all
invariant graphs of the RFDE.

\subsection{Finite mixed-jet prolongation}

For a parameter-dependent map \(f\), write
\begin{equation}
 J_{a,b,i,j}f
 =D_u^a\partial_\rho^b\partial_\nu^i\partial_\zeta^jf,
 \qquad n=a+b+i+j,
 \qquad p=b+i+j.                                     \label{eq:invdet-jet-notation}
\end{equation}
Let
\begin{equation}
 \mathcal I_k=\{(a,b,i,j)\in\mathbb N_0^4:
 b\leq3,\ i\leq2,\ j\leq2,\ a+b+i+j\leq k\}.     \label{eq:invdet-index-set}
\end{equation}
Order the indices first by increasing total grade \(n\) and, at equal
total grade, by increasing parameter grade \(p\).  Thus a derivative which
gains one state slot and loses one parameter derivative belongs to an
earlier block.  We bound the Picard jets through \(\mathcal I_{11}\) and
prove their convergence through \(\mathcal I_{10}\).  The last grade is
the single spatial reserve used when a highest derivative is evaluated on
two different flows.

For \(\eta=(a,b,i,j)\), the common Banach fiber is
\begin{equation}
 \mathfrak F_{\eta,N}
 =C\!\left(\Lambda_\delta;
 C_b\!\left(U;
 \mathcal L_{\rm sym}^a
   (U;U\times E_N)\right)\right),                    \label{eq:invdet-jet-fiber}
\end{equation}
with the supremum operator norm.  If several indices have the same
\((n,p)\), take their finite maximum product.  The fibers vary with \(N\),
but every estimate below uses their operator norms and has the same scalar
constant.

We spell out the flow induction.  With terminal time independent of the
parameters, a first parameter variation satisfies
\begin{equation}
 \frac{\dd}{\dd t}\partial_\lambda\Phi_Q^t
 =DQ(\Phi_Q^t)\partial_\lambda\Phi_Q^t
  +(\partial_\lambda Q)(\Phi_Q^t),
 \qquad \partial_\lambda\Phi_Q^0=0.                 \label{eq:invdet-first-flow-parameter}
\end{equation}
For \(a\geq2\), the pure state variation has the exact partition form
\begin{equation}
 \frac{\dd}{\dd t}D_u^a\Phi_Q^t
 =DQ(\Phi_Q^t)D_u^a\Phi_Q^t
 +\sum_{\substack{\mathcal P\text{ a partition of }\{1,\ldots,a\}\\
                   |\mathcal P|\geq2}}
 D^{|\mathcal P|}Q(\Phi_Q^t)
   \bigl[D_u^{|B|}\Phi_Q^t:B\in\mathcal P\bigr].   \label{eq:invdet-flow-partitions}
\end{equation}
After arbitrary mixed differentiation, the current flow jet occurs linearly
on the left.  The current jet of \(Q\) also occurs linearly.  If a parameter
derivative hits the argument of a vector-field jet, that vector-field jet
gains one state slot and loses one parameter derivative and is therefore in
an earlier block.  A product containing two positive-grade jets contains
only strictly lower total grades.  Finally, substituting
\(t=\vartheta-\rho r\) converts each derivative of the terminal time into a
factor \(r\) and a time derivative
\(\partial_t\Phi_Q^t=Q(\Phi_Q^t)\).  These formulas prove, by variation of
constants, the bound \cref{eq:invdet-tame-flow}; subtracting two such
systems uses the reserved state grade and gives its difference version.

Start Picard iteration at \(Z_N^{(0)}=(q_{0,\delta},0)\) and put
\(Z_N^{(k+1)}=\mathcal T_N(Z_N^{(k)})\).  Each finite iterate is smooth.
The following lemma is the precise triangular replacement for the phrase
``differentiate the same estimate.''

\begin{lemma}[Highest-block equation]
\label[lemma]{lem:invdet-highest-block}
Suppose all blocks preceding \((n,p)\subset\mathcal I_{11}\) range in
fixed bounded subsets of their common fibers.  Then the current block has
the form
\begin{equation}
 \mathbf J_{n,p}\mathcal T_N(Z)
 =\mathbf b_{n,p}(\mathbf J_{<n,p}Z)
  +\rho\mathscr L_{n,p}(\mathbf J_{<n,p}Z)
        [\mathbf J_{n,p}Z],                           \label{eq:invdet-highest-block-form}
\end{equation}
except that the first pure spatial block need only satisfy the corresponding
difference estimate.  There is one common tame bound
\begin{align}
 \|\rho\mathscr L_{n,p}\|&\leq\eta_\delta,
                                                               \label{eq:invdet-diagonal-small}\\
 \|\mathbf J_{n,p}\mathcal T_N(Z)
       -\mathbf J_{n,p}\mathcal T_N(\widetilde Z)\|
 &\leq \eta_\delta
       \|\mathbf J_{n,p}Z-\mathbf J_{n,p}\widetilde Z\|
       +C_{n,p,\delta}d_{<n,p}(Z,\widetilde Z),       \label{eq:invdet-highest-block-difference}
\end{align}
where
\begin{equation}
 \eta_\delta\leq C\delta\mathfrak B_\delta
 \int_0^\infty(1+r^{M_*})e^{-a_*r/2}\,\dd r=o(1)    \label{eq:invdet-jet-smallness}
\end{equation}
uniformly in \(N\).  The coefficient \(C_{n,p,\delta}\) is independent of
\(N\), though it may have the tame target dependence
\(P(R_\delta)e^{cR_\delta}\).
\end{lemma}

\begin{proof}
Apply the finite Fa\`a di Bruno formula to
\cref{eq:invdet-TQ,eq:invdet-TH}, using
\cref{eq:invdet-first-flow-parameter,eq:invdet-flow-partitions} for every
flow factor.  A factor carrying the full grade occurs at most once.  The
ordering in \cref{eq:invdet-index-set} places every same-total-grade term
created by differentiating a flow argument in an earlier block.  Terms with
two positive-grade factors use lower total grades.

More concretely, replace every occurrence of the current graph jet in these
formulas by an arbitrary element \(V\) of the current common fiber.  Retain
the terms linear in \(V\) in the variational equations
\cref{eq:invdet-first-flow-parameter,eq:invdet-flow-partitions}, solve them
with zero initial value, substitute the resulting flow jets and the stable
component of \(V\) in the terms linear in \(V\) in the Fa\`a di Bruno
expansion, and finally apply the stable semigroup integral.  This defines
\(\mathscr L_{n,p}V\).  Setting \(V=0\) defines
\(\mathbf b_{n,p}\).

Both components of the transform have an external factor \(\rho\).  If a
\(\rho\)-derivative strikes it, then
\begin{equation}
 \partial_\rho^b(\rho\mathcal F)
 =\rho\partial_\rho^b\mathcal F
   +b\partial_\rho^{b-1}\mathcal F,                  \label{eq:invdet-rho-Leibniz}
\end{equation}
and the second term belongs to a lower block.  Hence every current-block
term retains \(\rho\).  The tame-flow difference estimate and
\(\|e^{A_Nr}\|\leq e^{-a_*r}\) now give
\cref{eq:invdet-diagonal-small,eq:invdet-highest-block-difference}.  The
first pure spatial block is not affine because its first variational
equation contains \(DQ\) in the coefficient, but subtraction of the two
linear variational equations gives exactly the same difference estimate.
No coordinate sum over \(E_N\) occurs.
\end{proof}

\begin{proposition}[Existence and closure of the mixed jets]
\label[proposition]{prop:invdet-mixed-jets}
After decreasing \(\delta_0\), all Picard jets in \(\mathcal I_{11}\) lie
in invariant balls, and those in \(\mathcal I_{10}\) converge strongly in
their common fibers.  Their limits are the genuine mixed derivatives of
\cref{eq:invdet-fixed-graph}.  In particular, all rectangular derivatives
with
\begin{equation}
 0\leq a\leq3,\qquad 0\leq b\leq3,\qquad
 0\leq i\leq2,\qquad 0\leq j\leq2                \label{eq:invdet-required-rectangle}
\end{equation}
exist and are continuous.
\end{proposition}

\begin{proof}
Proceed through the finite ordered block list.  Once the radii of all lower
blocks have been fixed, \cref{eq:invdet-highest-block-form} gives
\begin{equation}
 \|\mathbf J_{n,p}Z_N^{(k+1)}\|
 \leq B_{n,p,\delta}
      +\eta_\delta\|\mathbf J_{n,p}Z_N^{(k)}\|.       \label{eq:invdet-invariant-jet-recurrence}
\end{equation}
Choose \(\delta_0\) so that the finitely many occurrences of
\(\eta_\delta\) are at most \(1/2\), and take the current radius larger than
\(2B_{n,p,\delta}+\|\mathbf J_{n,p}Z_N^{(0)}\|\).  This constructs invariant
balls through grade eleven, with radii independent of \(N\).

For consecutive iterates, suppose every preceding block converges
geometrically.  Equation \cref{eq:invdet-highest-block-difference} gives
\begin{equation}
 a_{k+1}\leq\eta_\delta a_k+Cq^k,
 \qquad \max\{\eta_\delta,q\}<1.                    \label{eq:invdet-jet-increment-recurrence}
\end{equation}
The increments are summable at every rate larger than
\(\max\{\eta_\delta,q\}\).  Induction proves strong convergence through
grade ten; grade eleven supplies the reserved derivative used in the
difference estimate.

It remains to identify the limits as derivatives.  If
\(f_k\to f\) and \(\partial_\lambda f_k\to g\) uniformly in a common fiber,
then on every line segment contained in the open parameter cylinder,
\begin{equation}
 f_k(\lambda+t)-f_k(\lambda)
 =\int_0^t\partial_\lambda f_k(\lambda+r)\,\dd r.     \label{eq:invdet-line-segment}
\end{equation}
Passing to the limit proves \(\partial_\lambda f=g\).  The same fundamental
theorem of calculus argument on state line segments identifies state
derivatives.  Repeating this in the block order proves that every limiting
jet is the derivative claimed in \cref{eq:invdet-required-rectangle}.
This is strong convergence in fixed Banach fibers; neither compactness nor
a weak-limit argument is used.
\end{proof}

For completeness, the inverse statement in \cref{lem:shared-graph-trace}
requires a declared norm on the whole prolongation.  Enumerate the finite
blocks by \(\ell=0,\ldots,L\).  Let \(C_{\ell k,\delta}\) bound the strict
lower-block coefficient from block \(k\) to block \(\ell\), and define
recursively
\begin{equation}
 w_{0,\delta}=1,\qquad
 w_{\ell,\delta}
 =1+4\sum_{k<\ell}C_{\ell k,\delta}w_{k,\delta},
 \qquad
 \|V\|_{\rm pr,\delta}
 =\max_{0\leq\ell\leq L}
     w_{\ell,\delta}^{-1}\|V_\ell\|.                \label{eq:invdet-prolonged-norm}
\end{equation}
The weights may depend tamely on the frozen target \(\delta\), but not on
\(N\).  Each diagonal residual block is
\(I-\rho\mathscr L_\ell\); hence, by
\cref{eq:invdet-diagonal-small},
\begin{equation}
 \|(I-\rho\mathscr L_\ell)^{-1}\|\leq2.             \label{eq:invdet-diagonal-inverse}
\end{equation}
Solving the lower-triangular residual system in order and using
\(\sum_{k<\ell}C_{\ell k,\delta}w_{k,\delta}/w_{\ell,\delta}\leq1/4\)
gives
\begin{equation}
 \|V\|_{\rm pr,\delta}\leq4\|F\|_{\rm pr,\delta}.  \label{eq:invdet-prolonged-inverse}
\end{equation}
Thus the full finite prolonged inverse is uniform in \(N\) and \(\delta\)
in the explicit scaled norm \cref{eq:invdet-prolonged-norm}.  An unweighted
product norm need not have a \(\delta\)-independent inverse bound, because
the strict lower-block coefficients can have the tame size
\(P(R_\delta)e^{cR_\delta}\).

Taylor's integral identity in the dummy amplitude is now legitimate.  Put
\begin{equation}
 \widehat Z_{j,N,\nu,\zeta}
 =\frac1{j!}\partial_\rho^jZ_{N,\rho,\nu,\zeta}\big|_{\rho=0},
 \qquad j=1,2.
\end{equation}
Since \(Z_{N,0,\nu,\zeta}=(q_{0,\delta},0)\),
\begin{equation}
 Z_{N,\delta,\nu,\zeta}
 =(q_{0,\delta},0)+\delta\widehat Z_{1,N,\nu,\zeta}
  +\delta^2\widehat Z_{2,N,\nu,\zeta}
  +\delta^3\widehat{\mathcal R}_{3,N},               \label{eq:invdet-Taylor}
\end{equation}
where
\begin{equation}
 \widehat{\mathcal R}_{3,N}
 =\frac12\int_0^1(1-t)^2
      \partial_\rho^3Z_{N,t\delta,\nu,\zeta}\,\dd t.             \label{eq:invdet-Taylor-remainder}
\end{equation}

We now define, rather than identify implicitly, the coefficients used in
the main text.  For a vector field \(\widehat V\) in the
\((\chi,d)\)-coordinates, set
\begin{equation}
 (\mathcal P_{\mathcal K}\widehat V)(X,Y)
 =D\mathcal K^{-1}(\mathcal K(X,Y))\,
        \widehat V(\mathcal K(X,Y)).                  \label{eq:invdet-vector-field-pullback}
\end{equation}
In particular, the reduced vector field and graph used in the main text are
\begin{equation}
 Q^{XY}_{N,\rho,\nu,\zeta}
 =\mathcal P_{\mathcal K}Q_{N,\rho,\nu,\zeta},
 \qquad
 H^{XY}_{N,\rho,\nu,\zeta}
 =H_{N,\rho,\nu,\zeta}\circ\mathcal K.                \label{eq:invdet-full-graph-pullback}
\end{equation}
Write
\[
 \widehat Z_{\ell,N}=(\widehat q_{\ell,N},\widehat H_{\ell,N}),
 \qquad
 \widehat{\mathcal R}_{3,N}
   =(\widehat{\mathcal R}^{\,Q}_{3,N},
      \widehat{\mathcal R}^{\,H}_{3,N}).
\]
The pulled-back coefficients are
\begin{align}
 q_{\ell,N}&=\mathcal P_{\mathcal K}\widehat q_{\ell,N},
 &H_{\ell,N}&=\widehat H_{\ell,N}\circ\mathcal K,
 \qquad \ell=1,2,\notag\\
 \mathcal R_{3,N}
 &=\mathcal P_{\mathcal K}\widehat{\mathcal R}^{\,Q}_{3,N},
 &\mathcal R^{\,H}_{3,N}
 &=\widehat{\mathcal R}^{\,H}_{3,N}\circ\mathcal K.   \label{eq:invdet-coefficient-pullbacks}
\end{align}
Because \(\mathcal K\) is independent of
\((\rho,\nu,\zeta)\), it commutes with all parameter differentiations and
with the integral Taylor formula.  On the retained set
\(q_{0,\delta}=q_0^{\chi d}\), and
\(\mathcal P_{\mathcal K}q_0^{\chi d}=q_0\).  Applying
\(\mathcal P_{\mathcal K}\) to the first component of
\cref{eq:invdet-Taylor} therefore gives the exact main-text expansion
\[
 Q^{XY}_{N,\delta,\nu,\zeta}
 =q_0+\delta q_{1,N}+\delta^2q_{2,N}
       +\delta^3\mathcal R_{3,N},
\]
which is \cref{eq:shared-graph-expansion}.  Since
\(\mathcal K,\mathcal K^{-1}\) are polynomial of degree at most two, the
state derivatives through the finite order used here lose only a
polynomial factor in \(R_\delta\); the weights
\(\langle s\rangle^Me^{c|s|}\) absorb this loss.

\begin{proposition}[Growing-tube coefficient envelope]
\label[proposition]{prop:invdet-coefficient-envelope}
Fix \(C_0,c_0>0\) and an integer \(M_0\geq0\).  There are
\(\delta_0>0\), \(C,c>0\), and an integer \(M\), depending on these three
fixed tube data and on the common preparation data but independent of \(N\),
such that, for every \(0<\delta<\delta_0\) and \(\ell=1,2\), the reduced
coefficients after conversion back to the \((X,Y)\)-coordinates satisfy
\begin{equation}
 \left|D_u^a\partial_\nu^i\partial_\zeta^j
 q_{\ell,N}(\gamma_0(s)+\xi,\nu,\zeta)\right|
 \leq C e^{c|s|}\langle s\rangle^M,                  \label{eq:invdet-coefficient-envelope}
\end{equation}
whenever \(0\leq a\leq3\), \(0\leq i,j\leq2\),
\(|s|\leq S_\delta+1\), \(\nu\in I_\nu\),
\(|\zeta|\leq\zeta_*\), and
\[
 |\xi|\leq C_0\delta e^{c_0|s|}\langle s\rangle^{M_0}.
\]
After increasing the output constants \(C,c,M\), without changing the
input tube, the same bound holds for the stable
coefficients \(H_{1,N},H_{2,N}\).  Moreover,
\begin{equation}
 \partial_\zeta q_{1,N}=0.                            \label{eq:invdet-q1-zeta-zero}
\end{equation}
\end{proposition}

\begin{proof}
Differentiate the fixed-point equations at \(\rho=0\).  Since
\(Z_{N,0,\nu,\zeta}=(q_{0,\delta},0)\), the first two reduced coefficients
in the \((\chi,d)\)-coordinates are, without a formal substitution,
\begin{align}
 \widehat q_{1,N}
 &=F_{N,\delta}(\mathcal E_{q_{0,\delta},0};
                       0,\nu,\zeta),\notag\\
 \widehat q_{2,N}
 &=\left.
   \left\{\partial_\rho F_{N,\delta}(\mathcal E_Z;\rho,\nu,\zeta)
   +D_{\mathcal E}F_{N,\delta}(\mathcal E_Z;\rho,\nu,\zeta)
          [\partial_\rho\mathcal E_Z]\right\}
   \right|_{\rho=0}.                                  \label{eq:invdet-q1-q2-formulas}
\end{align}
Here \(\partial_\rho\mathcal E_Z|_{\rho=0}\) is obtained from one
finite-time flow variation \cref{eq:invdet-first-flow-parameter} and
\(\widehat H_{1,N}\).  The stable formulas in the graph coordinates are
\begin{align}
 \widehat H_{1,N}(u)
 &=\int_0^\infty e^{A_Nr}
 G_{N,\delta}(\mathcal E_{q_{0,\delta},0}(u);
                         0,\nu,\zeta)\,\dd r,\notag\\
 \widehat H_{2,N}(u)
 &=\left.\partial_\rho
 \int_0^\infty e^{A_Nr}
 G_{N,\delta}\bigl(
   \mathcal E_Z(\Phi_Q^{-\rho r}u);
   \rho,\nu,\zeta\bigr)\,\dd r\right|_{\rho=0}.        \label{eq:invdet-H1-H2-formulas}
\end{align}
In the second line, differentiating the base shift contributes
\(-r q_{0,\delta}(u)\); every other new term contains \(Z_{1,N}\), one
fixed-time flow variation, or one derivative of the prepared data.

Equations \cref{eq:invdet-q1-q2-formulas,eq:invdet-H1-H2-formulas} show
that the first coefficient samples one fixed-delay backtrack and the second
samples at most a sum of two delays.  The factors \(r^k\), \(k=0,1\), in
the stable formulas are controlled by
\cref{eq:invdet-semigroup-moments}.  State and parameter differentiation
does not enlarge this delay depth: the atoms are fixed, and the
corresponding flow variations remain on the intervening finite singular
flow segments.

For \(|s|\leq S_\delta+1\), all of these points belong to the continuous
depth-two uncut hull in \cref{def:canonical-preparation}.  The raw
polynomials \cref{eq:invdet-raw-F,eq:invdet-raw-G}, the singular flow
variational equations, and \cref{eq:shared-products} therefore bound each
factor by \(P(\langle s\rangle)e^{c|s|}\), uniformly in \(N\).  The finite
Fa\`a di Bruno sums created by at most three state and two derivatives in
each parameter only multiply finitely many such weights; increasing
\(c\) and \(M\) absorbs every product.  The semigroup moments are uniform
by \cref{eq:invdet-semigroup-moments}.  The chain rule for
\(\mathcal K^{-1}\) introduces only further polynomial factors in
\(\langle s\rangle\), so the same envelope holds in \((X,Y)\).  This proves
\cref{eq:invdet-coefficient-envelope}.

Finally, before conjugation, at \(\rho=0\), the only possible direct
\(\zeta\)-dependence of the reduced \((X,Y)\)-formula is
\[
 K\pi_N^T\sum_{k=0}^mR_{k,N}\mathbf1(X_0-X_k),
\]
which vanishes because \(\pi_N^TR_{k,N}=0\) for every \(k\).
The base graph and singular flow are independent of \(\zeta\).  Thus the
unconjugated first coefficient has zero \(\zeta\)-derivative.  Since
\(\mathcal K\) is parameter independent, its push-forward and the pullback
in \cref{eq:invdet-vector-field-pullback} preserve this zero, proving
\cref{eq:invdet-q1-zeta-zero}.
\end{proof}

We record the local tame bound used in the main proof, including the norms
in which the block induction is performed.  Split the stable integral after
every derivative at \(L\log(1/\delta)\).  By choosing \(L\) after the finite
jet grade, its tail is bounded by
\begin{equation}
 C\mathfrak B_\delta
 \{1+L\log(1/\delta)\}^{M_*}
 e^{-a_*L\log(1/\delta)/2}=O(\delta^6).              \label{eq:invdet-convolution-tail}
\end{equation}
On the finite part, the additional base-flow displacement is at most
\(\delta L\log(1/\delta)=o(1)\).

Enumerate the finite ordered jet blocks by
\(\ell=0,\ldots,L_J\), including the order-zero block.  Fix any input-tube
data \(K_{\rm in},c_{\rm in}>0\) and \(m_{\rm in}\in\mathbb N_0\).
Choose fixed time buffers
\begin{equation}
 0=B_0<B_1<\cdots<B_{L_J},\qquad
 B_{j+1}-B_j>2\Theta_*+2,                            \label{eq:invdet-local-buffers}
\end{equation}
and, recursively, fixed tube weights
\[
 W_j^{\rm tub}(s)=e^{c_j^{\rm tub}|s|}
                   \langle s\rangle^{m_j^{\rm tub}},
 \qquad K_j\geq K_{\rm in},
\]
starting with
\(W_0^{\rm tub}=e^{c_{\rm in}|s|}\langle s\rangle^{m_{\rm in}}\)
and \(K_0=K_{\rm in}\).  The finite-time singular-flow estimates allow the
later parameters to be chosen monotonically,
\[
 c_{j+1}^{\rm tub}\geq c_j^{\rm tub},\qquad
 m_{j+1}^{\rm tub}\geq m_j^{\rm tub},\qquad
 K_{j+1}\geq K_j,
\]
and so that every current or fixed-delay evaluation and every finite
convolution shift used at level \(0\leq j<L_J\) maps
\begin{equation}
 \mathcal T^{\chi d}_{j,\delta}
 :=\left\{\mathcal K(\gamma_0(s)+\xi):
       |s|\leq S_\delta+1+B_j,
       |\xi|\leq K_j\delta W_j^{\rm tub}(s)\right\}   \label{eq:invdet-local-tubes}
\end{equation}
into \(\mathcal T^{\chi d}_{j+1,\delta}\).  Thus these are
\((\chi,d)\)-coordinate tubes, as required by the graph transform; their
preimages under \(\mathcal K\) are the displayed \((X,Y)\)-tubes on which
the final estimates are asserted.  The monotone choices and the increasing
buffers give
\(\mathcal T^{\chi d}_{j,\delta}\subset
  \mathcal T^{\chi d}_{j+1,\delta}\).
Only finitely many such choices are made.  Their planar normal radii remain
inside the shrinking agreement
strip for small \(\delta\), because
\(\delta e^{c_j^{\rm tub}S_\delta}
\langle S_\delta\rangle^{m_j^{\rm tub}}=o(\delta^\kappa)\).
The stable components remain inside their fixed saturation ball by the
global graph radii.  Planar bases outside the central interval enter only
fixed-width prepared endpoint strips; those are controlled by the global
tame majorant below.

For a block-valued map \(V\) and a positive weight \(W\), define
\begin{equation}
 \|V\|_{j;W}
 :=\sup_{\lambda\in\Lambda_\delta}
   \sup_{\substack{|s|\leq S_\delta+1+B_j\\
                   |\xi|\leq K_j\delta W_j^{\rm tub}(s)}}
       \frac{\|V(\mathcal K(\gamma_0(s)+\xi);\lambda)\|_{\rm op}}{W(s)}.
                                                               \label{eq:invdet-local-weighted-norm}
\end{equation}
Choose output weights
\(W_\ell^{\rm out}(s)=e^{c_\ell^{\rm out}|s|}
\langle s\rangle^{m_\ell^{\rm out}}\) recursively in the block order.
At the \(\ell\)-th step, enlarge its two exponents only enough to dominate
the finitely many products of earlier weights and fixed-shift ratios in the
source \(\mathbf b_\ell\) of
\cref{eq:invdet-highest-block-form}.  This is a finite, noncircular
construction: \(W_\ell^{\rm out}\) is chosen only after all weights occurring
in its source have been fixed.

On the inner portions \(|s|\leq S_\delta-C_{\rm buf}\), with
\(C_{\rm buf}\) chosen after the finite block list, all sampled points lie
in the uncut physical hull; the polynomial chart and its singular
variational equations give the corresponding output weight.  On the
complementary portions of the fixed buffered tubes,
\(\bigl||s|-S_\delta\bigr|\) is bounded, so the global tame majorant
\(P(R_\delta)e^{cR_\delta}\) is bounded by the chosen
\(CW_\ell^{\rm out}(s)\).
The global invariant radius of every block obtained from
\cref{eq:invdet-invariant-jet-recurrence} has the tame size
\[
 R_{\ell,\delta}^{\rm glob}\leq P_\ell(R_\delta)e^{c_\ell R_\delta}.
\]
This follows by the same finite induction because every source radius has
that form and \((1-\eta_\delta)^{-1}\leq2\).  After division by the
appropriate local weight, estimate the \(\ell\)-th block on
\(\mathcal T^{\chi d}_{L_J-\ell,\delta}\).  Every lower block is sampled at
worst on \(\mathcal T^{\chi d}_{L_J-\ell+1,\delta}\), which is contained in
the domain \(\mathcal T^{\chi d}_{L_J-k,\delta}\) already controlled for
\(k<\ell\).
The explicit highest-block formula
\cref{eq:invdet-highest-block-form} therefore gives
\begin{equation}
 \|\mathbf J_\ell Z\|_{L_J-\ell;W_\ell^{\rm out}}
 \leq B_\ell^{\rm loc}
   +\eta_\delta R_{\ell,\delta}^{\rm glob}
   +\sum_{k<\ell}C_{\ell k}
       \|\mathbf J_k Z\|_{L_J-k;W_k^{\rm out}},       \label{eq:invdet-local-block-recurrence}
\end{equation}
where \(B_\ell^{\rm loc}\) and \(C_{\ell k}\) are independent of
\(N,\delta\),
and
\[
 \eta_\delta R_{\ell,\delta}^{\rm glob}
 =\delta P(R_\delta)e^{cR_\delta}=o(1)
\]
after increasing the fixed polynomial and exponential constants.
The tail
\cref{eq:invdet-convolution-tail} is absorbed in \(B_\ell^{\rm loc}\).
Induction in \(\ell\), starting with the order-zero graph bound, gives one
uniform radius for every required block on its declared tube, hence in
particular on the original tube \(\mathcal T^{\chi d}_{0,\delta}\).
Applying this to
the integral Taylor remainder
\cref{eq:invdet-Taylor-remainder}, then applying the pullback
\cref{eq:invdet-coefficient-pullbacks} and absorbing its polynomial
coordinate loss, gives constants \(C_{\rm rem},c_{\rm rem}>0\) and
\(M_{\rm rem}\in\mathbb N_0\), depending on the fixed input-tube data but
not on \(N,\delta\), for which
\begin{equation}
 |D_u^a\partial_\nu^i\partial_\zeta^j
       \mathcal R_{3,N}(\gamma_0(s)+\xi)|
 \leq C_{\rm rem}\langle s\rangle^{M_{\rm rem}}
          e^{c_{\rm rem}|s|},
 \qquad a\leq3,\quad i,j\leq2,                       \label{eq:invdet-local-remainder}
\end{equation}
whenever
\[
 |s|\leq S_\delta+1,\qquad
 |\xi|\leq K_{\rm in}\delta e^{c_{\rm in}|s|}
                     \langle s\rangle^{m_{\rm in}}.
\]
This also makes explicit why a global
\(C_b^3\) remainder estimate independent of \(\delta\) is neither claimed
nor needed.

The same split gives the local graph size used in
\cref{lem:shared-graph-trace}.  Fix \(L,C_0>0\) and take
\(u=\gamma_0(s)+\xi\) in the \((X,Y)\)-coordinates, with
\(|s|\le L\) and \(|\xi|\le C_0\delta\).  The finite part of every stable
integrand and its required fixed-delay backtracks then remain in one fixed
uncut physical neighborhood and are bounded independently of \(\delta\);
their contribution is \(O(\delta)\).  The complementary tail is bounded by
\(\delta A_\delta O(\delta^6)=o(\delta)\).  Consequently
\begin{equation}
 \sup_{\substack{N,\nu,\zeta,\ |s|\le L\\|\xi|\le C_0\delta}}
 \|H_{N,\delta,\nu,\zeta}(\gamma_0(s)+\xi)\|_N
 \leq C_{L,C_0}\delta.                               \label{eq:invdet-fixed-compact-size}
\end{equation}

At \(\rho=0\), one differentiation of
\cref{eq:invdet-history-evaluation} samples at most one fixed-delay
backtrack.  A second differentiation samples at most a sum of two delays
and a variational integral along the intervening singular-flow segment.
For the stable component, the only additional terms are the moments
\begin{equation}
 \left\|\int_0^\infty r^ke^{A_Nr}\,\dd r\right\|
 \leq\frac{k!}{(D\gamma)^{k+1}},
 \qquad k=0,1.                                      \label{eq:invdet-semigroup-moments}
\end{equation}
Consequently, on the retained tube, the restrictions of \(Z_{1,N}\) and
\(Z_{2,N}\) depend only on the physical data on the declared depth-one and
depth-two hulls, respectively.  Their global prepared continuations need not
have this independence.  The full fixed point still uses the untruncated integral in
\cref{eq:invdet-TH}; the split in
\cref{eq:invdet-convolution-tail} is only an estimate.

\subsection{Exact and injective complete histories}

Fix the physical value \(\rho=\delta>0\) and suppress the parameters.  For
\(u_*\in U\), let
\begin{equation}
 u(s)=\Phi_Q^su_*,\qquad h(s)=H(u(s)),\qquad
 g(s)=G_{N,\delta}(\mathcal E_{Q,H}(u(s));
                         \delta,\nu,\zeta).
\end{equation}
Changing variables \(\tau=s-\delta r\) in the stable fixed-point equation
gives the exact identity
\begin{equation}
 h(s)=\int_{-\infty}^s
 e^{A_N(s-\tau)/\delta}g(\tau)\,\dd\tau.             \label{eq:invdet-stable-complete-convolution}
\end{equation}
For \(s\geq s_0\), split this integral at \(s_0\) and use the semigroup
law:
\begin{equation}
 h(s)=e^{A_N(s-s_0)/\delta}h(s_0)
 +\int_{s_0}^s e^{A_N(s-\tau)/\delta}g(\tau)\,\dd\tau.            \label{eq:invdet-mild-stable-equation}
\end{equation}
This is precisely the mild stable equation associated with
\cref{eq:invdet-raw-normal-form}; no differentiation of the convolution is
needed.  The first fixed-point equation supplies the classical reduced
equation, with the complete graph history
\cref{eq:invdet-history-lift}.  Uniqueness for the finite-dimensional
retarded equation therefore yields the semiconjugacy
\begin{equation}
 \mathcal S_{N,\nu,\zeta}^s\mathcal I_{Q,H}(u)
 =\mathcal I_{Q,H}(\Phi_Q^su),
 \qquad s\geq0,                                      \label{eq:invdet-semiconjugacy}
\end{equation}
for the prepared RFDE.

Define present-state restriction by
\begin{equation}
 \mathcal R_N(\Psi)=\pi_U\Psi(0).
\end{equation}
Then
\begin{equation}
 \mathcal R_N\mathcal I_{Q,H}(u)=u,\qquad
 D\mathcal R_ND\mathcal I_{Q,H}(u)=\Id_U.            \label{eq:invdet-left-inverse}
\end{equation}
Hence \(\mathcal I_{Q,H}\) is injective, a homeomorphism onto its image,
and, by \cref{prop:invdet-mixed-jets}, a \(C^3\) injective immersion.  Thus
it is a \(C^3\) embedding into
\(C([ -\Theta_*,0];X_N)\).  Equality of two lifted histories is equivalent
to equality of their reduced present states; it is not merely an equality
after stationary projection.

To return to the phase space \(\mathcal H_{N,\delta}\) in
\cref{sec:model-results}, first use the time reparametrization
\(\vartheta=\delta\theta_{\rm phys}\).  The voltage history determines
\begin{align*}
 X(\vartheta)
 &=\delta^{-1}\{\pi_N^Tv(\vartheta/\delta)-1\},\qquad
 -\Theta_*\le\vartheta\le0,\\
 h(\vartheta)
 &=\delta^{-2}P_{\perp,N}v(\vartheta/\delta)-g_NX(\vartheta)^2.
\end{align*}
The current resource value gives
\(Y(0)=\delta^{-2}(2/3-w(0))\), and the exact chart equation
\(Y'=-X+\delta\nu\) then recovers the discarded resource past by
\begin{equation}
 Y(\vartheta)=Y(0)+\int_\vartheta^0X(r)\,\dd r
                    +\delta\nu\,\vartheta .          \label{eq:invdet-recover-resource-history}
\end{equation}
These formulas define a continuous left inverse on the image of the slow
history lift.  Hence the map into
\(C([ -\Theta_*/\delta,0];\R^N)\times\R\) is injective even though the RFDE
phase space stores only the current resource value.  Consequently the
embedding and history-equality statements above concern the full original
RFDE state, not only its slow-coordinate representative.

Finally, let \(J\) be a retained interval of a graph orbit.  If
\begin{equation}
 \{\Phi_Q^{s+\vartheta}u_*:
          s\in J,\ -\Theta_*\leq\vartheta\leq0\}
\end{equation}
lies in the planar agreement region and all corresponding values of \(H\)
lie in the stable agreement ball \(\|h\|_N<h_*\), then every current and
delayed slot in the prepared right-hand side equals its unprepared
counterpart.  Equation \cref{eq:invdet-semiconjugacy} is therefore an exact
semiconjugacy for the physical RFDE on that retained segment.  In the
application to the selected traces, the identity
\[
 d(\gamma_0(s)+(U,V))=V+sU-\alpha U^2
\]
and the trace bound in \cref{eq:shared-trace-jets} give
\[
 |d|\leq C\delta e^{cS_\delta}\langle S_\delta\rangle^{M+1}
       =o(\delta^\kappa).
\]
The same bound holds after the finitely many retained delay backtracks.
Together with \(\|H\|_N\leq C\delta A_\delta=o(1)<h_*\), this verifies the
agreement criterion and proves the exact retained-history assertion.

\begin{lemma}[Precise relation with \cref{lem:shared-graph-trace}]
\label[lemma]{lem:invdet-to-main-lemma}
Assume
\cref{eq:shared-markov,eq:shared-curvature,eq:shared-layer-bounds} and fix a
preparation satisfying \cref{def:canonical-preparation}, including its
tame-flow clause \cref{def:invdet-tame-preparation}.  Then the following clauses of
\cref{lem:shared-graph-trace} follow from this subsection:
\begin{enumerate}[label=\textup{(\roman*)}]
 \item the existence, local uniqueness, invariance, and injectivity of the
 complete-history graph;
 \item the expansion \cref{eq:shared-graph-expansion}, the fixed-compact
 graph size, the coefficient envelope
 \cref{eq:invdet-coefficient-envelope}, and the local remainder estimate
 \cref{eq:shared-graph-remainder};
 \item the finite block-lower-triangular mixed prolongation through the
 derivative rectangle displayed in that lemma, with diagonal inverse bound
 two and full inverse bound four in the declared scaled norm
 \cref{eq:invdet-prolonged-norm};
 \item the exact-history assertion for every retained trace segment which
 satisfies the bootstrap furnished by \cref{eq:shared-trace-jets}.
\end{enumerate}
The one-sided trace existence, the endpoint normalization, the trace
linear isomorphisms, and the estimates in
\cref{eq:shared-trace-jets} are not proved here; they remain the separate
finite-dimensional trace argument beginning with
\cref{eq:shared-trace-operator}.

The explicit tame-flow clause excludes the uncontrolled fixed-delay flow
factor described in \cref{rem:invdet-missing-flow-hypothesis}.  Likewise, a uniform bound for
the full prolonged inverse is justified in the scaled norm
\cref{eq:invdet-prolonged-norm}, not in an undeclared unweighted product
norm.
\end{lemma}

\begin{proof}
The semigroup constants are \(1\) and \(D\gamma\) by
\cref{eq:poisson-bound}; the observation norm is at most \(m+2\) by
\cref{eq:shared-evaluation-TV}; and all nonlinear operator bounds are
dimension independent by \cref{eq:shared-products}.  Therefore
\cref{lem:invdet-C0-contraction} applies with constants independent of
\(N\).  The mixed-jet, coefficient, and Taylor assertions are
\cref{prop:invdet-mixed-jets,prop:invdet-coefficient-envelope,%
eq:invdet-Taylor,eq:invdet-local-remainder,%
eq:invdet-fixed-compact-size}, and the prolonged inverse is
\cref{eq:invdet-diagonal-inverse,eq:invdet-prolonged-inverse}.  Finally,
\cref{eq:invdet-mild-stable-equation,eq:invdet-left-inverse} give the exact
injective history lift, while the last bootstrap paragraph gives physical
agreement on the retained pieces.  These are precisely the four listed
clauses and no trace statement.
\end{proof}
